\begin{theorem}
Let $E$ be a finite set, let $F$ be a front on an infinite set $X$, and let
$f:F\to E$ be a super-sequence.  Then there are an infinite set $Y$, a
front $F'$ on $Y$ with $F'\subseteq F$, and a color $c\in E$ such that the
restriction of $f$ to $F'$ is constantly $c$.
\end{theorem}
\begin{proof}
We prove the assertion by induction on the cardinality of the finite
codomain.  The induction assertion is uniform in the front, its infinite
base, and the value map.  If the codomain is empty, there is no such value
map: by the density clause in \noderef{Front}, the front has a member, and
the map would have to assign that member a value in the empty set.  Thus this
case is vacuous.

Now suppose that $E$ is nonempty and choose a color $e\in E$.  Consider the
fiber
\[
   S_e=\{s\in F\mid f(s)=e\}.
\]
Apply \noderef{NashWilliams} to the subset $S_e\subseteq F$.  We obtain a
front $F_0$ on an infinite base $Y_0$, an inclusion $F_0\subseteq F$, and
exactly one of the following two alternatives:
\[
  \text{(selected color)}\quad F_0\subseteq S_e,
  \qquad\text{or}\qquad
  \text{(complement)}\quad F_0\cap S_e=\emptyset.
\]

In the selected-color case, use $F_0$, $Y_0$, its front structure, and its
inclusion into $F$.  The definition of a super-sequence in
\noderef{SuperSequence} lets us restrict the value map along that inclusion;
the definition of $S_e$ then gives
$f(s)=e$ for every $s\in F_0$.  This is the required constant
sub-super-sequence, with constant $e$.

In the complement case, every $s\in F_0$ satisfies $f(s)\ne e$.  Put
$E_0=E\setminus\{e\}$.  This is finite and has cardinality strictly smaller
than that of $E$.  Define a value map $g:F_0\to E_0$ by
$g(s)=f(s)$, regarding $f(s)$ as an element of $E_0$ using
$F_0\cap S_e=\emptyset$.  Together with the front $F_0$ on $Y_0$, this is
again a super-sequence in the sense of \noderef{SuperSequence}.  The
induction hypothesis supplies a front $F_1$ on an infinite base $Y_1$, an
inclusion $F_1\subseteq F_0$, and a color $c_0\in E_0$ such that
$g(s)=c_0$ for every $s\in F_1$.

Finally, compose $F_1\subseteq F_0$ with $F_0\subseteq F$.  The same front
$F_1$ is therefore a sub-front of the original domain, and the defining
equation $g(s)=f(s)$ shows that the original map $f$ is constantly the
underlying color $c_0\in E$ on $F_1$.  In either case the recursion removes a
color or terminates at the selected color, so the finite induction yields
the asserted constant sub-super-sequence.
\end{proof}
